\documentclass[12pt,oncampus]{article}
\usepackage{stfx_schedule}

\begin{document}

\stfxscheduleheader{Fall 2026}{RELS 398.10 / ENGL 398.10}{ST: Digital Methods for Historical Documents}

\begin{stfxschedule}

%% Week 1 -- setup week, no three-step cycle
\weekrow{Week 1}{Wed Sep 2 (\#1)\newline Fri Sep 4 (\#2)}{1. Orientation}{
  Finney, ``Manuscript Markup'' (pp.\ 323--360).\newline
  Griffitts, \S 3.4.1; TEI Guidelines vi.2.
}{
  \textit{Orientation week --- no three-step cycle.}\newline
  Local Git version control fundamentals (init, staging, atomic commits, diffs, commit log audit trail);
  Unicode polytonic Greek fonts and keyboard layouts; VS Code with XML tooling;
  GitHub, Transkribus and THREAD accounts.\newline
  \textbf{Choose your full-page manuscript witness by Fri Sep 4} (ungraded, but every later
  assignment builds on it).
}

%% Week 2
\weekrow{Week 2}{Tue Sep 8 (\#3)\newline Wed Sep 9 (\#4)\newline Fri Sep 11 (\#5)}{2. Cataloguing}{
  Griffitts \S 3.4.2, 3.6.1, 4.1; TEI Guidelines v.1--3.\newline
  Ceriani, \textit{Codex Marchalianus} introduction.
}{
  Manuscript identification, TEI Header, RDF triples and CTS URNs. XML and TEI
  fundamentals: elements, attributes, well-formedness, schema validation.\newline
  \textbf{Tue:} hand-write a \texttt{teiHeader} in a plain text editor.\newline
  \textbf{Wed:} validate it against the TEI schema in VS Code; build RDF triples.\newline
  \textbf{Fri:} AI-generated and enriched metadata, verified against your own header.\newline
  \textit{Add/Drop deadline: Mon Sep 14.}
}

%% Week 3
\weekrow{Week 3}{Tue Sep 15 (\#6)\newline Wed Sep 16 (\#7)\newline Fri Sep 18 (\#8)}{3. IIIF Imaging}{
  Griffitts \S 4.3--4.4.\newline
  Havens, ``Adobe Photoshop \& Digital Paleography''.
}{
  \textit{Infrastructure week --- the three-step cycle does not apply.}\newline
  International Image Interoperability Framework (IIIF), raking light,
  multispectral imaging, canvas linking; building a manifest for your passage.\newline
  \textbf{Cataloguing \& Metadata Assignment Due: Fri Sep 18}.
}

%% Week 4
\weekrow{Week 4}{Tue Sep 22 (\#9)\newline Wed Sep 23 (\#10)\newline Fri Sep 25 (\#11)}{4. Transcription\newline\textit{(Fri Sep 25: Remote Delivery)}}{
  Transkribus User Guide; THREAD Documentation.\newline
  Clemens \& Graham, \textit{Introduction to Manuscript Studies}.\newline
  stoa.org transcription conventions.
}{
  Diplomatic versus normalized transcription.\newline
  \textbf{Tue:} manual diplomatic transcription in a text editor --- this becomes
  your ground truth.\newline
  \textbf{Wed:} baseline detection and HTR model training in Transkribus.\newline
  \textbf{Fri (remote):} LLM-assisted transcription and correction.\newline
  \textbf{Transcription Assignment Due: Fri Sep 25} (manual diplomatic transcription).
}

%% Week 5 -- two meetings; the comparison IS the topic
\weekrow{Week 5}{Tue Sep 29 (\#12)\newline\textit{Wed Sep 30 Closed}\newline Fri Oct 2 (\#13)}{5. Reconciliation\newline\textit{(Wed Sep 30: Truth \& Reconciliation)}}{
  Haaf et al., ``Measuring Correctness of Double-Keying''.\newline
  Griffitts \S 4.5.2.
}{
  \textit{Two-step week by design: reconciliation is itself the comparison, so
  the missing Wednesday costs nothing.}\newline
  \textbf{Tue:} double-keying and consensus text building by hand.\newline
  \textbf{Fri:} character and word error rates for the HTR and AI passes, scored
  against your manual ground truth.\newline
  \textbf{Reconciliation Assignment Due: Fri Oct 2} (human vs.\ AI transcription).
}

%% Week 6
\weekrow{Week 6}{Tue Oct 6 (\#14)\newline Wed Oct 7 (\#15)\newline Fri Oct 9 (\#16)}{6. Normalization}{
  TEI Guidelines ch.\ 11 (Representation of Primary Sources).\newline
  Griffitts \S 4.6.1; Sabel, ``Encoding of Variant Taxonomies in TEI''.
}{
  Iotacism, orthographic variants, \texttt{tei:orig} / \texttt{tei:reg} markup.\newline
  \textbf{Tue:} regularize by hand, recording each decision.\newline
  \textbf{Wed:} rule-based normalization with a deterministic script.\newline
  \textbf{Fri:} AI-assisted normalization, checked against your own rules.\newline
  \textbf{Normalization Assignment Due: Fri Oct 9}.
}

%% Study Break
\shadedrow{Oct 12--16}{N/A}{\textbf{Fall Study Break} (No Classes)}{Thanksgiving Day Mon Oct 12.}

%% Week 7
\weekrow{Week 7}{Tue Oct 20 (\#17)\newline Wed Oct 21 (\#18)\newline Fri Oct 23 (\#19)}{7. Collation}{
  CollateX Documentation; Haentjens Dekker et al.\ (2015).\newline
  Griffitts \S 4.6.2--3, 4.7; Robinson, ``Collation and Textual Criticism''.
}{
  Variant taxonomy, base texts, alignment algorithms.\newline
  \textbf{Tue:} collate two witnesses by hand on paper.\newline
  \textbf{Wed:} automated alignment with the CollateX engine.\newline
  \textbf{Fri:} AI-assisted variant classification and adjudication.
}

%% Week 8
\weekrow{Week 8}{Tue Oct 27 (\#20)\newline Wed Oct 28 (\#21)\newline Fri Oct 30 (\#22)}{8. Apparatus}{
  TEI P5 Guidelines, Critical Apparatus module.\newline
  Burnard, \textit{What is the Text Encoding Initiative?}.
}{
  \texttt{tei:app}, \texttt{tei:lem}, \texttt{tei:rdg}; the TEI Critical
  Apparatus Toolbox.\newline
  \textbf{Tue:} encode an apparatus entry by hand.\newline
  \textbf{Wed:} generate the apparatus from CollateX output; validate.\newline
  \textbf{Fri:} AI-assisted apparatus drafting and error hunting.\newline
  \textbf{Textual Variants \& Collation Assignment Due: Fri Oct 30}.\newline
  \textit{Mid-term grades due Mon Oct 26.}
}

%% Week 9
\weekrow{Week 9}{Tue Nov 3 (\#23)\newline Wed Nov 4 (\#24)\newline Fri Nov 6 (\#25)}{9. Lemmatization}{
  CLTK (Classical Language Toolkit) Documentation.\newline
  TEI Guidelines ch.\ 18; Schmid, ``TreeTagger for Historical Languages''.
}{
  Lemmatization, morphological analysis, part-of-speech tagging.\newline
  \textbf{Tue:} parse and lemmatize a passage by hand with a lexicon.\newline
  \textbf{Wed:} CLTK and Morpheus tagging pipelines.\newline
  \textbf{Fri:} LLM-assisted parsing, scored against your manual parse.\newline
  \textit{Last day to drop: Wed Nov 4.}
}

%% Week 10 -- two meetings; deliberate slack
\weekrow{Week 10}{Tue Nov 10 (\#26)\newline\textit{Wed Nov 11 Closed}\newline Fri Nov 13 (\#27)}{10. Consolidation\newline\textit{(Wed Nov 11: Remembrance Day)}}{
  No new readings. Review your own edition files to date.
}{
  \textit{Consolidation week --- no new method, and none of the three steps.}\newline
  GitHub collaborative workflows (remotes, forks, pull requests, issue tracking);
  collaborative peer review of editions in progress; validation sweep of every TEI file built
  so far; catch-up on outstanding stages.\newline
  \textbf{Linguistic Markup Assignment Due: Fri Nov 13}.
}

%% Week 11
\weekrow{Week 11}{Tue Nov 17 (\#28)\newline Wed Nov 18 (\#29)\newline Fri Nov 20 (\#30)}{11. Alignment}{
  Leipzig Glossing Rules; TEI \S 16.1 (Link \& Alignment).\newline
  Tiedemann, ``Parallel Data Extraction''; Baumgardt, ``Digital Interlinear Editions''.
}{
  Sentence-level translation alignment and word-level interlinear glossing
  (Alpheios Alignment Editor, Logeion).\newline
  \textbf{Tue:} align and gloss by hand.\newline
  \textbf{Wed:} automated sentence alignment and glossing tools.\newline
  \textbf{Fri:} AI-assisted alignment and gloss generation, verified by hand.
}

%% Week 12
\weekrow{Week 12}{Tue Nov 24 (\#31)\newline Wed Nov 25 (\#32)\newline Fri Nov 27 (\#33)}{12. Commentary}{
  TEI \S 17 (Simple Analytic Mechanisms); Open Annotation Data Model.\newline
  Vander Meulen, ``Digital Commentary''.
}{
  Annotating passages, intertextual cross-referencing (Tesserae), commentary
  structures and linked scholarly apparatus.\newline
  \textbf{Tue:} write commentary notes by hand and encode them.\newline
  \textbf{Wed:} annotation tooling and cross-reference networks.\newline
  \textbf{Fri:} AI-assisted commentary drafting --- and its limits.\newline
  \textbf{Interlinear Glossing \& Translation Alignment Assignments Due: Fri Nov 27}.
}

%% Week 13
\weekrow{Week 13}{Tue Dec 1 (\#34)\newline Wed Dec 2 (\#35)\newline\textit{Thu Dec 3 Closed}\newline Fri Dec 4 (\#36)}{13. Publication}{
  Sahle, \textit{Digital Scholarly Editions} (vol.\ 2).\newline
  Pierazzo, \textit{Digital Textual Scholarship} (ch.\ 4--5).\newline
  TEI Publisher documentation.
}{
  \textit{Studio weeks --- no three-step cycle.} Stemma codicum, editorial
  principles, and writing the edition introduction; TEI validation, CSS
  rendering, web and print output.\newline
  \textbf{Scholarly Commentary Assignment Due: Fri Dec 4}.\newline
  \textit{Thu Dec 3: Closed (Feast Day). Fri Dec 4 runs a Wednesday timetable,
  10:30 AM --- the term's second-last class day.}
}

%% Exam Period
\eventrow{Dec 10--19}{N/A}{Comprehensive Digital Edition Portfolio}{
  \textbf{DECEMBER EXAM PERIOD}: submit the assembled TEI-conformant edition ---
  \texttt{teiHeader} with catalogue information, diplomatic transcription,
  normalized text with variants, lemmas and parsings, translation, and
  commentary.\newline
  \textbf{Edition Introduction, Print-Ready Publication, and Digital Edition
  Portfolio all due in this window.}
}

\end{stfxschedule}

\schedulefootnote{*Class meetings: Tuesdays 8:30 AM, Wednesdays 10:30 AM, Fridays 9:30 AM (MULH 3032). A standard week runs the same method three ways: \textbf{Tuesday} by hand in a text editor, \textbf{Wednesday} with specialist deterministic software, \textbf{Friday} with AI assistance (LLMs, Antigravity, Claude Code, or Hermes). Your Tuesday work is the ground truth against which Wednesday and Friday are measured. Weeks 1, 3, 10 and 13 depart from this pattern and say so in their row.}

\end{document}
